% Replace the two consecutive statements maps:src:F.3.6 and night:P1.8,
% together with both proofs, by this single statement and proof.
% Place night:P1.15 and Definition def:hb before this statement so that the
% algebraic splitting, contraction and h_b are already defined.
\begin{theorem}[Lifted periods and their exact parameter domain]
\hypertarget{tag-02O8}{}\hypertarget{tag-02QZ}{}
\label{maps:src:F.3.6}\label{night:P1.8}
\StatementRecord{night:P1.8}{proved}{02QZ}{\SeriesRef{night:P1.5}, \SeriesRef{night:P1.15}}
Let $P$ be the rank-three wedge bundle of \cref{night:P1.5}, and use
its normalized contraction $\alpha$ and the algebraic splitting $\sigma_b$.
The planes
\begin{equation}
 F_\xi=(\sigma_b+\xi\alpha e)(P),\qquad \xi\in\mathbb C,
 \hypertarget{tag-02ES}{}\label{maps:eq:F.21}
\end{equation}
are holomorphic rank-three subbundles of the marked rank-six system.
They extend over its finite root charts of orders five and four.
The function $h_b$ of \cref{def:hb} has an attained finite minimum
$h_{b,\min}$, tends to $+\infty$ at the cusp, and satisfies
\begin{equation}
 F_\xi\cap\overline{F_\xi}=0
 \quad\Longleftrightarrow\quad
 h_b+\operatorname{Im}\xi\ne0
 \hypertarget{tag-02ET}{}\label{maps:eq:F.22}
\end{equation}
at each nonsingular base point, including the finite root points.
Thus the parameters transverse on the completed base are exactly
$\operatorname{Im}\xi>-h_{b,\min}$.
In an integral marking at the cusp their periods are
\begin{equation}
 Z(q)=\frac{\log q}{2\pi i}I_3+C_\xi(q),
 \qquad C_\xi\text{ holomorphic at }q=0.
 \hypertarget{tag-02EJ}{}\label{maps:eq:F.13}
\end{equation}
\end{theorem}
\begin{proof}
On the root curve of orders five and eight, let $\mathcal H_D$ be
the logarithmic de Rham bundle of the pencil. Its Hodge subbundle is
$\mathcal E=T^3\oplus T$, with $\ell=T^3$. The polarization identifies
$\mathcal H_D/\mathcal E$ with $\mathcal E^\vee$. This remains true
at the node: the dualizing differentials have independent limits,
and their isotropic pairing extends there holomorphically.
The primitive projection $p$ and the marked extension give
\begin{align}
 0&\longrightarrow\mathbb C e\longrightarrow\mathcal V
   \longrightarrow\mathcal W\longrightarrow0,
 \hypertarget{tag-02EK}{}\label{maps:eq:F.14}\\
 P&=p(\ell\wedge\mathcal H_D).
 \hypertarget{tag-02EL}{}\label{maps:eq:F.15}
\end{align}
In logarithmic frames the first sequence extends over the cusp,
since the nilpotent exponential commutes with its inclusion and
quotient. It extends equivariantly over the finite root charts.
The wedge map is injective after primitive projection, as proved
in \cref{night:P1.5}. The Hodge sequence therefore gives
\begin{equation}
 0\longrightarrow T^4\longrightarrow P
   \longrightarrow\mathcal O\oplus T^2\longrightarrow0.
 \hypertarget{tag-02EM}{}\label{maps:eq:F.16}
\end{equation}
For the coarse projection $\rho$, invariant sections in the two
root charts give
\begin{equation}
 \rho_*T^j=\mathcal O_{\mathbb P^1}
 \left(\left\lfloor\frac{2j}{5}\right\rfloor+
       \left\lfloor-\frac{3j}{8}\right\rfloor\right).
 \hypertarget{tag-02EN}{}\label{maps:eq:F.17}
\end{equation}
Finite-group averaging makes descent exact. For $j=\pm2,\pm4$
the degree is $-1$, so both cohomology groups vanish. Consequently
$P=\mathcal O\oplus T^2\oplus T^4$ and
$H^1(P^\vee)=0$. The pullback of \eqref{maps:eq:F.14} to $P$
splits holomorphically. Its splittings differ by $z\alpha e$,
since $H^0(P^\vee)=\mathbb C\alpha$. We use the algebraic choice
$\sigma_b$ and its marked contraction from \cref{night:P1.15}.

The real-intersection calculation and well-definedness proof of
\cref{def:hb} give $P\cap\bar P=\mathbb Cn$ with
$\alpha(n)=1$ and
$\sigma_b(n)-\overline{\sigma_b(n)}=2ih_b e$.
The lifted vectors for parameter $\xi$ differ from their
conjugates by $2i(h_b+\operatorname{Im}\xi)e$.
An arbitrary complex multiple of $n$ gives no further
intersection direction. This proves \eqref{maps:eq:F.22}.

We calculate its cusp behavior in the fixed companion marking.
For a local nonzero $u\in\ell_0$, put $r=ip(u\wedge\bar u)$
and $A=\alpha(r)$, so that $n=r/A$.
Put $D=C_fC_g^{-1}-I$, $N_V=(BA)^{-1}-I$ and let $N_W$ be its
quotient. The only nonzero column of $D$ is
$(0,1,1,1)^t$, and the top row of $N_V$ is
$(0,0,0,-1,1,-1)$. The wedge projection and normalized contraction
satisfy, for all $a,b\in\mathbb C^4$,
\[
 (N_V)_{0,*}(0,p(a\wedge Db))
 =\operatorname{contr}(a\wedge Db)=a_0b_0,
 \qquad N_Wp(a\wedge Db)=0.
\]
These identities follow by expanding the two-by-two minors in
the companion basis. Write $u=\exp(sD)a(q)$ with
$s=\log q/(2\pi i)$ and $y=\operatorname{Im}s$.
The vector $a$ is holomorphic, and $a_0(0)\ne0$ because
$dx/y$ has nonzero nodal residue. In the logarithmic frame,
\[
 e^{-sN_W}r=ip(a\wedge\bar a)+2y\,p(a\wedge D\bar a),
 \qquad
 A=i\operatorname{contr}(a\wedge\bar a)+2y|a_0|^2.
\]
The splitting is bounded in that frame. On returning by
$e^{sN_V}$, its imaginary invariant-coordinate contribution is
$2|a_0|^2y^2+O(y)$. Division by
$A=2|a_0|^2y+O(1)$ yields
$h_b(q)=y+O(1)$, uniformly with $\operatorname{Re}s$ in a unit
interval. Thus $h_b(q)\to+\infty$ at the cusp.

The same identities give the logarithmic graph. At the limiting
plane choose a symplectic basis $a,b,d_1,d_2$ with $N_Hb=a$
and $\ell=\mathbb Cb$. Then
\begin{equation}
 P_0=\langle p(b\wedge a),b\wedge d_1,b\wedge d_2\rangle,
 \quad N_W(b\wedge d_i)=a\wedge d_i,
 \quad N_Wp(b\wedge a)=0.
 \hypertarget{tag-02EU}{}\label{maps:eq:F.23}
\end{equation}
The first generator maps under $N_V$ to the nonzero invariant
line, while the other two map to independent quotient directions.
Therefore $N_V:F_{\xi,0}\to\ker N_V$ is an isomorphism for
every $\xi$. In the primitive integral Jordan marking, this
makes the untwisted plane a graph with holomorphic coefficient
$C_\xi(q)$. Restoring $e^{sN_V}$ adds $sI_3$, giving
\eqref{maps:eq:F.13}.

The height extends continuously over both finite root charts.
Its cusp growth and compactness away from that cusp give an
attained minimum. The completed punctured base is connected, so
its image is $[h_{b,\min},\infty)$. Equation
\eqref{maps:eq:F.22} gives the exact parameter half-plane.
The fourth power of the order-eight action on $\mathcal H_D$
is $-I$; it acts trivially on $\mathcal W,P,\alpha$ and
$\mathcal V$. The equivariant lifted plane therefore descends
to the order-four chart. The other chart has order five.
\end{proof}
